% Module 4 section file. Included by ../module4_alfven_continuum_eigenmodes.tex.
\section{Shear-Alfv\'en waves in a uniform plasma}
\label{sec:uniform_alfven}

\subsection{Geometry and simplifying assumptions}
Consider a uniform equilibrium
\begin{equation}
\mathbf B_0=B_0\widehat{\mathbf z},
\qquad
\rho_0=\text{constant},
\qquad
p_0=\text{constant},
\qquad
\mathbf J_0=\mathbf 0.
\end{equation}
The unit vector $\widehat{\mathbf z}$ points along the equilibrium magnetic field. Take a transverse displacement
\begin{equation}
\boldsymbol\xi=\xi_x(z,t)\widehat{\mathbf x},
\end{equation}
with no $x$ or $y$ dependence. This displacement is perpendicular to $\mathbf B_0$, and
\begin{equation}
\nabla\cdot\boldsymbol\xi=0.
\end{equation}
Therefore, Eq.~\eqref{eq:p1_xi} gives $p_1=0$. The perturbation is shear-like rather than compressional.

\subsection{Magnetic perturbation produced by field-line bending}
Using Eq.~\eqref{eq:b1_xi},
\begin{align}
\mathbf B_1
&=\nabla\times(\boldsymbol\xi\times\mathbf B_0)\\
&=\nabla\times\left[-B_0\xi_x(z,t)\widehat{\mathbf y}\right]\\
&=B_0\frac{\partial\xi_x}{\partial z}\widehat{\mathbf x}.
\label{eq:b1_uniform_alfven}
\end{align}
A spatially uniform sideways translation, $\partial_z\xi_x=0$, does not bend the field and creates no restoring force. A displacement that varies along $z$ tilts neighboring field-line segments differently and produces a transverse magnetic perturbation.

The perturbed current is
\begin{align}
\mathbf J_1
&=\frac{1}{\mu_0}\nabla\times\mathbf B_1
=\frac{B_0}{\mu_0}\frac{\partial^2\xi_x}{\partial z^2}\widehat{\mathbf y}.
\end{align}
The Lorentz force is then
\begin{align}
\mathbf J_1\times\mathbf B_0
&=\frac{B_0^2}{\mu_0}\frac{\partial^2\xi_x}{\partial z^2}
\widehat{\mathbf x}.
\end{align}
For a sinusoid, the second derivative has the opposite sign to the displacement, so this is a restoring force.

\subsection{Wave equation and Alfv\'en speed}
The $x$ component of the linear momentum equation becomes
\begin{equation}
\rho_0\frac{\partial^2\xi_x}{\partial t^2}
=\frac{B_0^2}{\mu_0}\frac{\partial^2\xi_x}{\partial z^2}.
\end{equation}
Define the Alfv\'en speed
\begin{equation}
\boxed{v_A=\frac{B_0}{\sqrt{\mu_0\rho_0}}.}
\label{eq:alfven_speed}
\end{equation}
The wave equation is
\begin{equation}
\boxed{
\frac{\partial^2\xi_x}{\partial t^2}
=v_A^2\frac{\partial^2\xi_x}{\partial z^2}.}
\label{eq:shear_alfven_wave_equation}
\end{equation}
Thus, magnetic tension plays the role that string tension plays for a stretched string, while mass density supplies inertia.

\subsection{Plane-wave dispersion relation}
Insert
\begin{equation}
\xi_x=\widetilde\xi_x e^{i(k_\parallel z-\omega t)}.
\end{equation}
Then $\partial_t^2\rightarrow-\omega^2$ and $\partial_z^2\rightarrow-k_\parallel^2$, giving
\begin{equation}
\boxed{\omega^2=k_\parallel^2v_A^2,\qquad \omega=\pm k_\parallel v_A.}
\label{eq:uniform_alfven_dispersion}
\end{equation}
The two signs represent propagation in opposite directions along the magnetic field. When reporting a positive spectrum, one normally uses $\omega=|k_\parallel|v_A$.

The phase and group speeds are
\begin{equation}
v_{\mathrm{ph},\parallel}=\frac{\omega}{k_\parallel}=\pm v_A,
\qquad
v_{\mathrm g,\parallel}=\frac{\partial\omega}{\partial k_\parallel}=\pm v_A.
\end{equation}
In ideal uniform MHD the branch is nondispersive: all parallel wavelengths propagate at the same speed. Geometry, compressibility, kinetic effects, and finite-Larmor-radius physics introduce dispersion in more complete models.

\subsection{Polarization}
For the shear-Alfv\'en branch,
\begin{equation}
\boldsymbol\xi\cdot\mathbf B_0=0,
\qquad
\mathbf B_1\cdot\mathbf B_0=0,
\qquad
\nabla\cdot\boldsymbol\xi\approx0.
\end{equation}
The displacement and magnetic perturbation are transverse, while propagation is primarily parallel. This distinguishes the branch from a compressional Alfv\'en or fast magnetosonic wave, which changes magnetic-field magnitude and couples more strongly to density and pressure compression.

\subsection{Energy partition}
The perturbed fluid velocity is $\mathbf v_1=\partial_t\boldsymbol\xi$. The cycle-averaged kinetic and magnetic energy densities for a plane wave are
\begin{align}
\langle\mathcal K\rangle
&=\frac{1}{4}\rho_0|\widetilde v_x|^2,\\
\langle\mathcal M\rangle
&=\frac{1}{4\mu_0}|\widetilde B_{1x}|^2.
\end{align}
From Eq.~\eqref{eq:b1_uniform_alfven}, $|\widetilde B_{1x}|=B_0|k_\parallel\widetilde\xi_x|$, while $|\widetilde v_x|=|\omega\widetilde\xi_x|$. Using Eq.~\eqref{eq:uniform_alfven_dispersion},
\begin{equation}
\frac{|\widetilde B_{1x}|^2}{\mu_0}
=\rho_0|\widetilde v_x|^2.
\end{equation}
Therefore,
\begin{equation}
\boxed{\langle\mathcal K\rangle=\langle\mathcal M\rangle.}
\end{equation}
The wave oscillates between fluid kinetic energy and magnetic bending energy, with equal cycle averages in the ideal uniform limit.

\subsection{Elsasser variables}
For transverse incompressible perturbations, define velocity-unit magnetic perturbation
\begin{equation}
\mathbf b_v=\frac{\mathbf B_1}{\sqrt{\mu_0\rho_0}}
\end{equation}
and Elsasser variables
\begin{equation}
\mathbf z^\pm=\mathbf v_1\mp\mathbf b_v.
\end{equation}
The linear equations reduce to
\begin{equation}
\left(\frac{\partial}{\partial t}\pm v_A\frac{\partial}{\partial z}\right)\mathbf z^\pm=0.
\end{equation}
Thus $\mathbf z^+$ and $\mathbf z^-$ represent oppositely propagating Alfv\'en wave packets. This form is useful in turbulence theory, although the present eigenmode benchmark works with standing global modes rather than freely propagating packets.

\subsection{Characteristic magnitude}
For a deuterium--tritium plasma with $B_0=\SI{2.5}{T}$ and $\rho_0=3.3\times10^{-8}\ \mathrm{kg\,m^{-3}}$,
\begin{equation}
v_A\approx\frac{2.5}{\sqrt{(4\pi\times10^{-7})(3.3\times10^{-8})}}
\approx1.23\times10^7\ \mathrm{m\,s^{-1}}.
\end{equation}
For a parallel scale of order $R_0=\SI{3}{m}$, the angular-frequency scale is
\begin{equation}
\frac{v_A}{R_0}\approx4.09\times10^6\ \mathrm{s^{-1}},
\end{equation}
corresponding to $v_A/(2\pi R_0)\approx\SI{651}{kHz}$. This estimate explains why the benchmark frequencies are in the hundreds of kilohertz.

\subsection{Limits of the uniform derivation}
The result $\omega=|k_\parallel|v_A$ remains the organizing local relation, but a confined plasma is not uniform. Both $k_\parallel$ and $v_A$ vary with magnetic surface and Fourier harmonic. Toroidal geometry also couples harmonics that are independent in the slab. The next sections build those effects one at a time.
